% 04_numerical_method.tex
%
% Purpose:
%   Collect the numerical approximation, implementation components, diagnostics,
%   and computational workflow in one paper-style section.
%
% Main algorithm:
%   Summarizes Nyström/Kress discretization, proxy/MFS quasiperiodic Green
%   functions, density-to-Rayleigh extractors, mode-selection diagnostics, and
%   the end-to-end assembly workflow.
%
% Based on:
%   sections/discretization.tex, sections/truncation_and_diagnostics.tex, and
%   sections/numerical_workflow.tex from the first report draft.
%
% Main changes:
%   Moves diagnostics and workflow under the numerical method section rather
%   than leaving them as separate top-level sections.
%
% Numerical goal:
%   Identify the numerical parameters and checks needed before experiments are
%   considered reliable.

\section{Numerical Method and Implementation}

\subsection{Boundary integral discretization and Kress quadrature}

The inclusion boundaries are assumed smooth and are discretized by a periodic
Nystr\"om method.  For a \(2\pi\)-periodic boundary parameter \(t\), the nodes are
\[
  t_j=\frac{2\pi j}{N},
  \qquad
  j=0,\ldots,N-1,
\]
with even \(N\).  Smooth integral blocks use the periodic trapezoid rule, which
is spectrally accurate for smooth periodic data.

Logarithmic singularities are handled by Kress product quadrature
\cite{Kress1991,Kress2014}.  The kernel is split as
\[
  k(t,s)
  =
  \phi(t,s)\log\!\left(4\sin^2\frac{t-s}{2}\right)
  +
  \psi(t,s),
\]
where \(\phi\) and \(\psi\) are smooth.  The Kress weights for the singular
factor are
\[
  R_j^{(N/2)}(t)
  =
  -\frac{4\pi}{N}
  \left[
  \sum_{n=1}^{N/2-1}\frac{1}{n}\cos(n(t_j-t))
  +
  \frac{1}{N}\cos\!\left(\frac{N}{2}(t_j-t)\right)
  \right].
\]
The corresponding Nystr\"om entries have the form
\[
  A_{ij}=R_{ij}\phi(t_i,t_j)+h\psi(t_i,t_j),
  \qquad
  h=\frac{2\pi}{N}.
\]
For the Muller matrix, hypersingular differences such as
\[
  T^{(k_e)}-T^{(k_i)}
\]
are treated in a cancellation-aware Kress form.  The leading singular parts
cancel to a lower-order structure, but the implementation still needs the
correct split and diagonal treatment.

\subsection{Quasiperiodic Green functions and proxy periodization}

The quasiperiodic Green function is one of the delicate numerical parts of the
method.  In a one-direction periodic open waveguide, away from thresholds, a
representative Rayleigh expansion is
\[
  G_{\mathrm{QP}}^{(k)}((x,y),(x_s,y_s))
  =
  \frac{\ii}{2d}
  \sum_{m\in\mathbb Z}
  \frac{1}{\gamma_m}
  \ee^{\ii\beta_m(y-y_s)}
  \ee^{\ii\gamma_m|x-x_s|}.
\]
Direct lattice sums can be slowly convergent or singular near exceptional
parameters.  The codebase therefore uses MFS/proxy ideas to approximate
quasiperiodic corrections.  A typical proxy representation has the form
\[
  G_{\mathrm{QP}}^{(k)}(x,y)
  \approx
  G^{(k)}(x,y)
  +
  \sum_{\ell=1}^{P}q_\ell(y)\,G^{(k)}(x,z_\ell),
\]
where \(G^{(k)}\) is the free-space singular Green function and the proxy
sources \(z_\ell\) lie outside the physical cell.  The free-space singular part
is retained; the proxy correction is a homogeneous field inside the cell.  This
strategy is flexible, but its source geometry, check points, and truncation
parameters require convergence checks.  It should be understood alongside other
quasiperiodic integral representations such as those discussed by Barnett and
Greengard \cite{BarnettGreengard2010}.

\subsection{Far-field extractors}

The matrices \(F_\pm\) map the center boundary density to outgoing Rayleigh
coefficients at the ports:
\[
  s^- = F_-\eta,
  \qquad
  s^+ = F_+\eta .
\]
They are not additional unknowns.  At the continuous level, the coefficient
functions come from the Rayleigh expansion of \(G_{\mathrm{QP}}\).  For example,
for a source point \(z=(z_x,z_y)\) and right reference wall \(X_R\),
\[
  g_m^R(z)
  =
  \frac{\ii}{2\sqrt d\,\gamma_m}
  \exp[-\ii\beta_m z_y-\ii\gamma_m(z_x-X_R)] .
\]
The double-layer source-normal contribution has coefficient
\[
  -\ii(\gamma_m\nu_x+\beta_m\nu_y)g_m^R(z),
\]
and the single-layer contribution is proportional to \(-g_m^R(z)\) under the
current \(\eta=[\tau;-\sigma]\) convention.  The left extractor is analogous
with the left-going Rayleigh phase and the corresponding source-normal
derivative formula.

\subsection{Mode selection, truncation, and diagnostics}

For real \(k\), homogeneous Rayleigh orders satisfy
\[
  \gamma_m^2=k^2-\beta_m^2 .
\]
Thus the Rayleigh order \(m\) decays in an open homogeneous direction exactly
when
\[
  k<|\beta_m|.
\]
In the simple one-dimensional prototype with \(0\le\beta\le\pi/d\),
\[
  \text{all Rayleigh orders decay}
  \iff
  k<\min_m|\beta_m|=\beta .
\]
Periodic leads require Bloch multiplier selection instead; admissibility is
selected through
\[
  \Asc(k,\beta)z=\lambda\Bsc(k,\beta)z,
\]
not through a scalar Rayleigh light-cone formula.

When \(k\) is real and propagating Bloch modes are present, some multipliers may
lie on or near the unit circle.  Magnitude alone then does not decide incoming
versus outgoing.  A typical flux diagnostic is
\[
  \mathcal F_x
  =
  \operatorname{Im}\int_0^d \overline{u(y)}\,\partial_xu(y)\,\dd y
  \approx
  \operatorname{Im}\sum_m \overline{D_m}N_m .
\]
In early experiments it is often more robust to work in a band gap, or to use
\(k+\ii\varepsilon\) with \(\varepsilon>0\), so that unit-circle ambiguity is
moved away from the selection threshold.

The Rayleigh truncation parameter \(M\) controls trace resolution and the amount
of evanescent coupling retained between separated objects and ports.  For large
\(|m|\), real \(k\), and fixed \(\beta\),
\[
  \gamma_m\approx \ii|\beta_m|,
  \qquad
  |\beta_m|\sim\frac{2\pi|m|}{d}.
\]
An evanescent component crossing an open distance \(\ell\) has the heuristic
decay factor
\[
  \ee^{-\operatorname{Im}\gamma_m\ell}.
\]
If each inclusion is separated from its nearest cell wall by at least
\(\delta\), then the gap between adjacent inclusions across a port is roughly at
least \(2\delta\), giving the practical coupling heuristic
\[
  \ee^{-2\operatorname{Im}\gamma_m\delta}.
\]
This is a useful guideline for choosing \(M\), but it is not a rigorous global
error estimate.

Changing \(M\) changes the truncated scattering matrix and therefore the Bloch
generalized eigenproblem.  Convergence should be checked not only for Rayleigh
tails but also for multipliers, outgoing subspaces, singular values, and fields.
Useful diagnostics include:
\begin{itemize}
  \item the selected counts \(K_{\mathrm{out}}^\pm\), which should be near \(K\)
  in non-degenerate band-gap or regularized cases;
  \item the conditioning and rank of
  \[
    \Phi_{\mathrm{all}}^\pm
    =
    \begin{bmatrix}
    D_{\mathrm{in}}^\pm & D_{\mathrm{out}}^\pm\\
    N_{\mathrm{in}}^\pm & N_{\mathrm{out}}^\pm
    \end{bmatrix};
  \]
  \item principal angles or projector differences between \(E_{\mathrm{out}}\)
  at successive \(M\) values;
  \item convergence of \(\sigmin(\Atr(k,\beta))\);
  \item BIE row residuals, port matching residuals, and field-plot convergence.
\end{itemize}

\subsection{Implementation components}

The present code organization separates the numerical roles into packages and
driver scripts:
\[
\begin{array}{ll}
\texttt{+geom} & \text{geometry construction and cached contour data},\\
\texttt{+kernel} & \text{quasiperiodic Green, MFS/proxy, and Kress kernel data},\\
\texttt{+op} & \text{Muller matrix assembly, including }\texttt{construct\_A\_QP},\\
\texttt{+bloch} & \text{Rayleigh channels, scattering matrices, modes, traces, selection},\\
\texttt{+viz} & \text{geometry and field visualization helpers}.
\end{array}
\]
The scattering workflows are represented by scripts such as
\path|scat_ld_lead_in.m|, \path|scat_edc_lead_in.m|, and
\path|scat_edc_ps.m|.  Eigenvalue or transmission-eigenvalue scans are
represented by scripts with the \path|tep_| prefix.  These names are
implementation anchors rather than a fixed API specification.

\subsection{Computational workflow}

A typical computation proceeds as follows.
\begin{enumerate}
  \item Choose the geometry, transverse quasiperiodicity \(\beta\), and either a
  fixed \(k\) or a scan range for \(k\).
  \item Build the boundary discretization for the center obstacle and, when
  needed, the representative lead obstacles.
  \item Construct the quasiperiodic Green or proxy/MFS data required by the
  exterior boundary integral operators.
  \item Assemble the center Muller matrix \(\Aqp\) and the center homogeneous
  trace matrices \(H_\Sigma^a,H_\Sigma^b\).
  \item Build the Rayleigh channel data
  \[
    \beta_m,\qquad \gamma_m,\qquad E=\diag(\ee^{\ii\gamma_mL_0}).
  \]
  \item Assemble each lead-cell scattering matrix \(S_{\mathrm{cell}}^\pm\) and
  solve \(\Asc z=\lambda\Bsc z\).
  \item Convert Bloch eigenvectors to wall Cauchy traces and select incoming and
  outgoing half-lead trace subspaces.
  \item Build the density far-field extractors \(F_-\) and \(F_+\), and assemble
  the global trace-matching matrix \(\Atr\).
  \item For forced scattering, choose valid incoming coefficients
  \(p_{\mathrm{in}}^\pm\) and solve \(\Atr x=b_{\mathrm{in}}\).
  \item For an eigenvalue or TEP-type scan, compute
  \[
    \sigmin(\Atr(k,\beta))
  \]
  across the parameter range and refine candidate dips.
  \item Validate convergence in \(N_{\mathrm{tot}}\), \(M\), proxy parameters,
  and mode-selection thresholds.
\end{enumerate}
